\chapter{A Pipeline Where the People Could Answer}
\label{sec:11}
Andry José Hernández Romero is a makeup artist from Venezuela. On 29 August 2024 he presented himself at the San Ysidro port of entry to seek asylum \autocite{LULAC2025HernandezRomero}. He was detained while his claim was processed, passed his credible-fear interview, and had, his lawyer said, a strong case \autocite{Vega2025HernandezRomero}. He was held at the Otay Mesa Detention Center, which a private company runs under contract with ICE \autocite{LULAC2025HernandezRomero}. There, a company employee screening him flagged two tattoos, crowns above the words for “mom” and “dad,” as marks of the Tren de Aragua gang. The employee was a former police officer, and Hernández Romero's lawyers said the government's allegation of gang membership rested on the tattoos and on no other information \autocite{Palomera2025Blade}. On 14 March 2025 the President signed a proclamation invoking the Alien Enemies Act against members of that gang \autocite{procl10903}. The next day three flights left Texas for El Salvador, and Hernández Romero was on one of them \autocite{LULAC2025HernandezRomero,jgg2025contempt}. He had a hearing on his asylum claim scheduled for later that month. He didn't appear at it, because he was in El Salvador's Terrorism Confinement Center, CECOT, and for weeks his lawyer didn't know where he was. Asked by reporters for the evidence against him, DHS cited state secrets and ongoing litigation \autocite{Vega2025HernandezRomero}. He spent 125 days in CECOT and was released in a prisoner exchange that sent more than 250 Venezuelans to Venezuela, not back to the country where his claim was pending. His asylum application in Spain is pending \autocite{Wiggins2026Spain}.

Nearly every position where a refusal can be defeated appears in that account. A contractor's employee wrote the case description from two tattoos. The determination was a gang-membership designation, not a finding about him reached at any hearing. He had come directly and presented himself, the case Article 31 of the Refugee Convention protects, and his claim to protection was pending, the case the first entry of this chapter's schedule protects. The handoff was to a foreign government's prison, under an arrangement for which the United States paid El Salvador's government about \$6 million \autocite{Renteria2025ElSalvador}. On the evening of the flights, a federal judge ordered any plane carrying people removed under the proclamation returned to the United States, and the flights went on \autocite{jgg2025contempt}. When the remedy came, it was diplomatic. He was released to Venezuela, not returned to the proceeding he had been removed from.

This chapter runs the floor's procedure on an abduction and forcible transfer pipeline, where the people acted on could answer; Part IV runs it on a targeting pipeline, where they can't. Hernández Romero could have answered: he had a lawyer, a pending claim and a hearing date. What he lacked was time. The operator moved faster than notice. The Supreme Court held that people designated under the proclamation must receive notice “within a reasonable time and in such a manner as will allow them to actually seek habeas relief” before removal \autocite{jgg2025}, and later that about a day's notice, with no word of how to contest the removal, was not enough \autocite{aarp2025}. Both rulings came after he was in CECOT. A pipeline where the people can answer is one where the law makes the operator wait for the answer.

So this chapter's instrument is a statute. The agency that runs the pipeline is statutory, and so are its powers, its means and its money, and what one Congress granted another can withdraw. At the chapter's centre is the statute that would end the agency and the seven terms it must carry, written over the acts the agency performs rather than over the agency's name. The sections before it supply what the statute needs: the pipeline as it was assembled, the acts the statute's schedule forbids, and the positions where those acts move when one route is closed. The sections after it say what a floor can do before the statute passes. Those measures are interim and lapse when the statute takes effect. Amnesty International documented the agency's operations in four cities, and its U.S. section's executive director concluded that “ICE is fundamentally broken and cannot be fixed” \autocite{Amnesty2026PressRelease}. This chapter agrees, and asks what ending it has to carry.
